\textbf{第一关代码}
\begin{minted}{c++}
#include<bits/stdc++.h>
using namespace std;
int dp[31][31][410][610]; // i, j, k, l, 时间，地点，水量，食物。
int daystate[35] = {0, 2, 2, 1, 3, 1, 2, 3, 1, 2, 2,
     3, 2, 1, 2, 2, 2, 3, 3, 2, 2, 
     1, 1, 2, 1, 3, 2, 1, 1, 2, 2}; // 1, 2, 3晴朗，高温，沙暴；
int wa[5] = {0, 5, 8, 10};
int fo[5] = {0, 7, 6, 10};
vector<int> v[35];
void sta(){

    v[27].push_back(26); 
    v[27].push_back(21);
    
    v[26].push_back(27); 
    v[26].push_back(25);
    
    v[25].push_back(24); 
    v[25].push_back(1); 
    v[25].push_back(26);
    
    v[24].push_back(25); 
    v[24].push_back(23);
    
    v[23].push_back(24); 
    v[23].push_back(21);
    
    v[21].push_back(23); 
    v[21].push_back(9); 
    v[21].push_back(27);
    
    v[9].push_back(21); 
    v[9].push_back(10); 
    v[9].push_back(15);
    
    v[10].push_back(9); 
    v[10].push_back(11);
    
    v[11].push_back(10); 
    v[11].push_back(12);
    
    v[12].push_back(13); 
    v[12].push_back(11);
    
    v[13].push_back(12); 
    v[13].push_back(15);
    
    v[15].push_back(9); 
    v[15].push_back(13);
    
    v[1].push_back(25);
}

void dfs(int t, int cur, int w, int f){
    if(t == 0){
        printf("决策：移动，当天的天气%d 第几天：%d %d %d %d %d\n", daystate[t], t, cur, w, f, dp[t][cur][w][f]);
        return;
    }
    if(daystate[t] != 3){
        if(cur != 15){ //移动
            for(auto p : v[cur])
                if(dp[t - 1][p][w + 2 * wa[ daystate[t] ]][f + 2 * fo[ daystate[t] ]] == dp[t][cur][w][f]){
                    dfs(t - 1, p, w + 2 * wa[ daystate[t]], f + 2 * fo[ daystate[t]]);
                    printf("决策：移动，当天的天气%d 第几天：%d %d %d %d %d\n", daystate[t], t, cur, w, f, dp[t][cur][w][f]);
                    return;
                }
        }else {
            for(auto p : v[cur])
                for(int k = 2 * wa[daystate[t]]; k <= w + 2 * wa[daystate[t]]; k++) // 要走一步
                    for(int l = 2 * fo[daystate[t]]; l <= f + 2 * fo[daystate[t]]; l++)
                        if(dp[t - 1][p][k][l] == dp[t][cur][w][f] + 
                            10 * (w + 2 * wa[daystate[t]] - k) + 20 * (f + 2 * fo[daystate[t]] - l)){
                                dfs(t - 1, p, k, l);
                                printf("决策：买东西，当天的天气%d 第几天：%d %d %d %d %d\n", daystate[t], t, cur, w, f, dp[t][cur][w][f]);
                                return;
                            }
        }
    }
    if(cur == 12)
        if(dp[t - 1][cur][w + 3 * wa[daystate[t]]][f + 3 * fo[daystate[t]]] == dp[t][cur][w][f] - 1000){
            dfs(t - 1, cur, w + 3 * wa[daystate[t]] , f + 3 * fo[daystate[t]]);
            printf("决策：挖矿，当天的天气%d 第几天：%d %d %d %d %d\n", daystate[t], t, cur, w, f, dp[t][cur][w][f]);
            return;
        }
    
    if(dp[t - 1][cur][w + 1 * wa[daystate[t]]][f + 1 * fo[daystate[t]]] == dp[t][cur][w][f]){
        dfs(t - 1, cur, w + 1 * wa[daystate[t]] , f + 1 * fo[daystate[t]]);
        printf("决策：停留，当天的天气%d 第几天：%d %d %d %d %d\n", daystate[t], t, cur, w, f, dp[t][cur][w][f]);
        return;
    }
    
}
int main(){
    sta();

    for(int i = 0; i <= 30; i++)
        for(int j = 0; j <= 30; j++)
            for(int k = 0; k <= 400; k++)
                for(int l = 0; l <= 600; l++)
                    dp[i][j][k][l] = -5e3;

    for(int k = 0; k <= 400; k++){
        for(int l = 0; l <= (1200 - 3 * k) / 2; l++){
            if(10000 - 5 * k - 10 * l >= 0)
                dp[0][1][k][l] = 10000 - 5 * k - 10 * l;
            
        }
    }

    for(int i = 1; i <= 30; i++){
        if(daystate[i] != 3) //非沙暴天
            for(int j = 1; j <= 27; j++)
                for(int k = wa[daystate[i]] * 2; k <= 400; k++)   
                    for(int l = fo[daystate[i]] * 2; l <= (1200 - 3 * k) / 2; l++){
                        if(dp[i - 1][j][k][l] < 0) continue;
                        for(auto t : v[j])
                            dp[i][t][k - wa[daystate[i]] * 2][l - fo[daystate[i]] * 2] = 
                            max(dp[i - 1][j][k][l], dp[i][t][k - wa[daystate[i]] * 2][l - fo[daystate[i]] * 2]);
                            //移动
                    }
        
        for(int j = 1; j <= 27; j++)
            for(int k = wa[daystate[i]]; k <= 400; k++)   
                for(int l = fo[daystate[i]]; l <= (1200 - 3 * k) / 2; l++){
                    if(dp[i - 1][j][k][l] < 0) continue;
                    dp[i][j][k - wa[daystate[i]]][l - fo[daystate[i]]] = 
                    max(dp[i - 1][j][k][l], dp[i][j][k - wa[daystate[i]]][l - fo[daystate[i]]]);
                            //停留
                }

        for(int k = wa[daystate[i]] * 3; k <= 400; k++)   
            for(int l = fo[daystate[i]] * 3; l <= (1200 - 3 * k) / 2; l++){
                if(dp[i - 1][12][k][l] < 0) continue;
                dp[i][12][k - wa[daystate[i]] * 3][l - fo[daystate[i]] * 3] = 
                        max(dp[i - 1][12][k][l] + 1000, dp[i][12][k - wa[daystate[i]] * 3][l - fo[daystate[i]] * 3]);
                        //挖矿
            }

        for(int k = 0; k <= 400; k++)   
            for(int l = 0; l <= (1200 - 3 * k) / 2; l++){
                if(k > 0)
                    dp[i][15][k][l] = max(dp[i][15][k][l], dp[i][15][k - 1][l] - 10); //价格双倍
                if(l > 0)
                    dp[i][15][k][l] = max(dp[i][15][k][l], dp[i][15][k][l - 1] - 20);
            }//买东西
    }
    int res = 0;
    int t = 0;
    for(int i = 1; i <= 30; i++)
        if(res < dp[i][27][0][0]){
            res = dp[i][27][0][0];
            t = i;
        }
    dfs(t, 27, 0, 0);
    cout << res << " " << t << endl;
}
\end{minted}

\textbf{第二关代码}
\begin{minted}{c++}
#pragma GCC optimize("O2")
#pragma GCC optimize("Ofast")
#pragma GCC optimize("inline")
#pragma GCC optimize("omit-frame-pointer")
#pragma GCC optimize("unroll-loops")
#include<bits/stdc++.h>
using namespace std;
int Po = 27; 
int dp[31][65][403][603]; // i, j, k, l, 时间，地点，水量，食物。
int daystate[35] = {0, 2, 2, 1, 3, 1, 2, 3, 1, 2, 2,
     3, 2, 1, 2, 2, 2, 3, 3, 2, 2, 
     1, 1, 2, 1, 3, 2, 1, 1, 2, 2}; // 1, 2, 3晴朗，高温，沙暴；
int wa[5] = {0, 5, 8, 10};
int fo[5] = {0, 7, 6, 10};
vector<int> v[70];

void sta(){
    //第二关数据
    for(int i = 1; i <= 63; i++)
        if(i % 8){
            v[i].push_back(i + 1);
            v[i + 1].push_back(i);
        }
    
    for(int i = 1; i <= 56; i++){
            v[i].push_back(i + 8);
            v[i + 8].push_back(i);
        }
    
    for(int i = 0; i <= 3; i++)
        for(int j = 2; j <= 8; j++){
            v[i * 16 + j].push_back(i * 16 + j + 7);
            v[i * 16 + j + 7].push_back(i * 16 + j);
        }
    for(int i = 1; i <= 3; i++)
        for(int j = 1; j <= 7; j++){
            v[i * 16 - (8 - j)].push_back(i * 16 - (8 - j) + 9);
            v[i * 16 - (8 - j) + 9].push_back(i * 16 - (8 - j));
        }

}


void dfs(int t, int cur, int w, int f){
    if(t == 0){
        printf("决策：移动，当天的天气%d 第几天：%d %d %d %d %d\n", daystate[t], t, cur, w, f, dp[t][cur][w][f]);
        return;
    }
    if(daystate[t] != 3){
        if(cur != 39 && cur != 62){ //移动
            for(auto p : v[cur])
                if(dp[t - 1][p][w + 2 * wa[ daystate[t] ]][f + 2 * fo[ daystate[t] ]] == dp[t][cur][w][f]){
                    dfs(t - 1, p, w + 2 * wa[ daystate[t]], f + 2 * fo[ daystate[t]]);
                    printf("决策：移动，当天的天气%d 第几天：%d %d %d %d %d\n", daystate[t], t, cur, w, f, dp[t][cur][w][f]);
                    return;
                }
        }else{
            for(auto p : v[cur])
                for(int k = 2 * wa[daystate[t]]; k <= w + 2 * wa[daystate[t]]; k++) // 要走一步
                    for(int l = 2 * fo[daystate[t]]; l <= f + 2 * fo[daystate[t]]; l++)
                        if(dp[t - 1][p][k][l] == dp[t][cur][w][f] + 
                            10 * (w + 2 * wa[daystate[t]] - k) + 20 * (f + 2 * fo[daystate[t]] - l)){
                                dfs(t - 1, p, k, l);
                                printf("决策：买东西，当天的天气%d 第几天：%d %d %d %d %d\n", daystate[t], t, cur, w, f, dp[t][cur][w][f]);
                                return;
                            }
                        }
    }
    if(cur == 30 || cur == 55)
        if(dp[t - 1][cur][w + 3 * wa[daystate[t]]][f + 3 * fo[daystate[t]]] == dp[t][cur][w][f] - 1000){
            dfs(t - 1, cur, w + 3 * wa[daystate[t]] , f + 3 * fo[daystate[t]]);
            printf("决策：挖矿，当天的天气%d 第几天：%d %d %d %d %d\n", daystate[t], t, cur, w, f, dp[t][cur][w][f]);
            return;
        }

    if(dp[t - 1][cur][w + 1 * wa[daystate[t]]][f + 1 * fo[daystate[t]]] == dp[t][cur][w][f]){
        dfs(t - 1, cur, w + 1 * wa[daystate[t]] , f + 1 * fo[daystate[t]]);
        printf("决策：停留，当天的天气%d 第几天：%d %d %d %d %d\n", daystate[t], t, cur, w, f, dp[t][cur][w][f]);
        return;
    }
    
}

int main(){
    ios::sync_with_stdio(false);
    cin.tie(0);
    sta();
    
    for(int i = 0; i <= 30; i++)
        for(int j = 0; j <= 64; j++)
            for(int k = 0; k <= 400; k++)
                for(int l = 0; l <= 600; l++)
                    dp[i][j][k][l] = -5e3;

    for(int k = 0; k <= 400; k++){
        for(int l = 0; l <= (1200 - 3 * k) / 2; l++){
            if(10000 - 5 * k - 10 * l >= 0)
                dp[0][1][k][l] = 10000 - 5 * k - 10 * l;
            
        }
    }

    for(int i = 1; i <= 30; i++){
        if(daystate[i] != 3) //非沙暴天
            for(int j = 1; j <= 64; j++)
                for(int k = wa[daystate[i]] * 2; k <= 400; k++)   
                    for(int l = fo[daystate[i]] * 2; l <= (1200 - 3 * k) / 2; l++){
                        if(dp[i - 1][j][k][l] < 0) continue;
                        for(auto t : v[j])
                            dp[i][t][k - wa[daystate[i]] * 2][l - fo[daystate[i]] * 2] = 
                            max(dp[i - 1][j][k][l], dp[i][t][k - wa[daystate[i]] * 2][l - fo[daystate[i]] * 2]);
                            //移动
                    }
        
        for(int j = 1; j <= 64; j++)
            for(int k = wa[daystate[i]]; k <= 400; k++)   
                for(int l = fo[daystate[i]]; l <= (1200 - 3 * k) / 2; l++){
                    if(dp[i - 1][j][k][l] < 0) continue;
                    dp[i][j][k - wa[daystate[i]]][l - fo[daystate[i]]] = 
                    max(dp[i - 1][j][k][l], dp[i][j][k - wa[daystate[i]]][l - fo[daystate[i]]]);
                            //停留
                }

        for(int k = wa[daystate[i]] * 3; k <= 400; k++)   
            for(int l = fo[daystate[i]] * 3; l <= (1200 - 3 * k) / 2; l++){
                if(dp[i - 1][30][k][l] < 0) continue;
                dp[i][30][k - wa[daystate[i]] * 3][l - fo[daystate[i]] * 3] = 
                        max(dp[i - 1][30][k][l] + 1000, dp[i][12][k - wa[daystate[i]] * 3][l - fo[daystate[i]] * 3]);
                        //挖矿
            }
        
        for(int k = wa[daystate[i]] * 3; k <= 400; k++)   
            for(int l = fo[daystate[i]] * 3; l <= (1200 - 3 * k) / 2; l++){
                if(dp[i - 1][55][k][l] < 0) continue;
                dp[i][55][k - wa[daystate[i]] * 3][l - fo[daystate[i]] * 3] = 
                        max(dp[i - 1][55][k][l] + 1000, dp[i][55][k - wa[daystate[i]] * 3][l - fo[daystate[i]] * 3]);
                        //挖矿
            }

        for(int k = 0; k <= 400; k++)   
            for(int l = 0; l <= (1200 - 3 * k) / 2; l++){
                if(k > 0)
                    dp[i][39][k][l] = max(dp[i][39][k][l], dp[i][39][k - 1][l] - 10); //价格双倍
                if(l > 0)
                    dp[i][39][k][l] = max(dp[i][39][k][l], dp[i][39][k][l - 1] - 20);
            }//买东西
        for(int k = 0; k <= 400; k++)   
            for(int l = 0; l <= (1200 - 3 * k) / 2; l++){
                if(k > 0)
                    dp[i][62][k][l] = max(dp[i][62][k][l], dp[i][62][k - 1][l] - 10); //价格双倍
                if(l > 0)
                    dp[i][62][k][l] = max(dp[i][62][k][l], dp[i][62][k][l - 1] - 20);
            }//买东西
    }
    int res = 0;
    int t = 0;
    for(int i = 1; i <= 30; i++)
        if(res < dp[i][64][0][0]){
            res = dp[i][64][0][0];
            t = i;
        }

    dfs(t, 64, 0, 0);

    cout << res << " " << t << endl;

}
\end{minted}

\textbf{第三关代码}
\begin{minted}{c++}
#pragma GCC optimize("O2")
#pragma GCC optimize("Ofast")
#pragma GCC optimize("inline")
#pragma GCC optimize("omit-frame-pointer")
#pragma GCC optimize("unroll-loops")
#include<bits/stdc++.h>
using namespace std;
int dp[20][20][403][603]; // i, j, k, l, 时间，地点，水量，食物。
int daystate[35] = {0, 2, 2, 2, 2, 2, 2, 2, 2, 2, 2}; // 1, 2, 3晴朗，高温，沙暴；
int wa[5] = {0, 3, 9, 10};
int fo[5] = {0, 4, 9, 10};
vector<int> v[30];
int comes = 1, wt, fd;
void sta(){

    // 第三关数据
    v[1].push_back(2), v[1].push_back(4), v[1].push_back(5);
    v[2].push_back(1), v[2].push_back(3), v[2].push_back(4);
    v[3].push_back(2), v[3].push_back(4), v[3].push_back(8), v[3].push_back(9);
    v[4].push_back(1), v[4].push_back(2), v[4].push_back(3), 
    v[4].push_back(5), v[4].push_back(6), v[4].push_back(7);
    v[5].push_back(1), v[5].push_back(4), v[5].push_back(6);
    v[6].push_back(4), v[6].push_back(5), v[6].push_back(7),
    v[6].push_back(12), v[6].push_back(13);
    v[7].push_back(4), v[7].push_back(11), v[7].push_back(6), v[7].push_back(12);
    v[8].push_back(3), v[8].push_back(9);
    v[9].push_back(3), v[9].push_back(8), v[9].push_back(10), v[9].push_back(11);
    v[10].push_back(9), v[10].push_back(11), v[10].push_back(13);
    v[11].push_back(7), v[11].push_back(9), v[11].push_back(10),
    v[11].push_back(12), v[11].push_back(13);
    v[12].push_back(6), v[12].push_back(7), v[12].push_back(11), v[12].push_back(13);
    v[13].push_back(6), v[13].push_back(10), v[13].push_back(11), v[13].push_back(12);

    
}

struct S{
    int t, day, wat, foo, re;
}st[40];
int cnt = 0;

void dfs(int t, int cur, int w, int f, int now){
    if(now == t){
        wt = w;
        fd = f;
        comes = cur;
        st[cnt++] = {t, cur, w, f, dp[t][cur][w][f]};
        return;
    }

    if(t == 0){
        st[cnt++] = {t, cur, w, f, dp[t][cur][w][f]};
        // printf("%d %d %d %d %d\n", t, cur, w, f, dp[t][cur][w][f]);
        return;
    }
    if(daystate[t] != 3){
        if(cur != 20 && cur != 20){ //移动
            for(auto p : v[cur])
                if(dp[t - 1][p][w + 2 * wa[ daystate[t] ]][f + 2 * fo[ daystate[t] ]] == dp[t][cur][w][f]){
                    dfs(t - 1, p, w + 2 * wa[ daystate[t]], f + 2 * fo[ daystate[t]], now);
                    // printf("%d %d %d %d %d\n", t, cur, w, f, dp[t][cur][w][f]);
                    return;
                }
        }else{
            for(auto p : v[cur])
                for(int k = 2 * wa[daystate[t]]; k <= w + 2 * wa[daystate[t]]; k++) // 要走一步
                    for(int l = 2 * fo[daystate[t]]; l <= f + 2 * fo[daystate[t]]; l++)
                        if(dp[t - 1][p][k][l] == dp[t][cur][w][f] + 
                            10 * (w + 2 * wa[daystate[t]] - k) + 20 * (f + 2 * fo[daystate[t]] - l)){
                                dfs(t - 1, p, k, l, now);
                                // printf("%d %d %d %d %d\n", t, cur, w, f, dp[t][cur][w][f]);
                                return;
                            }
                        }
    }
    if(cur == 9)
        if(dp[t - 1][cur][w + 3 * wa[daystate[t]]][f + 3 * fo[daystate[t]]] == dp[t][cur][w][f] - 200){
            dfs(t - 1, cur, w + 3 * wa[daystate[t]] , f + 3 * fo[daystate[t]], now);
            // printf("%d %d %d %d %d\n", t, cur, w, f, dp[t][cur][w][f]);
            return;
        }

    if(dp[t - 1][cur][w + 1 * wa[daystate[t]]][f + 1 * fo[daystate[t]]] == dp[t][cur][w][f]){
        dfs(t - 1, cur, w + 1 * wa[daystate[t]] , f + 1 * fo[daystate[t]], now);
        // printf("%d %d %d %d %d\n", t, cur, w, f, dp[t][cur][w][f]);
        return;
    }
    
}

int main(){
    ios::sync_with_stdio(false);
    cin.tie(0);
    auto start = std::chrono::high_resolution_clock::now();
    sta();
    

    double trans_prob[3][3];
    trans_prob[1][1] = 0.5;   // 晴→晴: 2次
    trans_prob[1][2] = 0.5;   // 晴→高: 6次
    trans_prob[1][3] = 0.0;        // 晴→沙: 0
        
        // 高温(2)的转移（只考虑晴朗和高温）
    trans_prob[2][1] = 0.2;  // 高→晴: 5次
    trans_prob[2][2] = 0.8;  // 高→高: 10次
    trans_prob[2][3] = 0.0;        // 高→沙: 0
        
        // 沙暴(3)的转移（实际不会用到）
    trans_prob[3][1] = 0.0;
    trans_prob[3][2] = 0.0;
    trans_prob[3][3] = 0.0;
    srand(time(0));
    
    for(int day = 1; day <= 10; day++){

    double ran = rand() % 50000 / (double)50000;
    if(daystate[day - 1] == 0){
        if(trans_prob[1][1] >= ran) {
            
            daystate[day] = 1; // 晴天
        }else 
            daystate[day] = 2;// 高温
    }else {
        if(trans_prob[daystate[day - 1]][1] >= ran) // 根据上一天的天气来看
            daystate[day] = 1; // 晴天
        else 
            daystate[day] = 2;// 高温
    }
    // daystate[day] = 2;
    for(int i = 0; i <= 10; i++)
        for(int j = 0; j <= 13; j++)
            for(int k = 0; k <= 400; k++)
                for(int l = 0; l <= 600; l++)
                    if(i != day - 1 || comes != j || day == 1 || wt != k || fd != l)
                        dp[i][j][k][l] = -5e3;
    if(day == 1){
        for(int k = 0; k <= 400; k++){
            for(int l = 0; l <= (1200 - 3 * k) / 2; l++){
                if(10000 - 5 * k - 10 * l >= 0)
                    dp[0][1][k][l] = 10000 - 5 * k - 10 * l;
                
            }
        }
    }

    for(int i = 1; i <= 10; i++){
        if(daystate[i] != 3) //非沙暴天
            for(int j = 1; j <= 13; j++)
                for(int k = wa[daystate[i]] * 2; k <= 400; k++)   
                    for(int l = fo[daystate[i]] * 2; l <= (1200 - 3 * k) / 2; l++){
                        if(dp[i - 1][j][k][l] < 0) continue;
                        for(auto t : v[j])
                            dp[i][t][k - wa[daystate[i]] * 2][l - fo[daystate[i]] * 2] = 
                            max(dp[i - 1][j][k][l], dp[i][t][k - wa[daystate[i]] * 2][l - fo[daystate[i]] * 2]);
                            //移动
                    }
        
        for(int j = 1; j <= 13; j++)
            for(int k = wa[daystate[i]]; k <= 400; k++)   
                for(int l = fo[daystate[i]]; l <= (1200 - 3 * k) / 2; l++){
                    if(dp[i - 1][j][k][l] < 0) continue;
                    dp[i][j][k - wa[daystate[i]]][l - fo[daystate[i]]] = 
                    max(dp[i - 1][j][k][l], dp[i][j][k - wa[daystate[i]]][l - fo[daystate[i]]]);
                            //停留
                }
        
        for(int k = wa[daystate[i]] * 3; k <= 400; k++)   
            for(int l = fo[daystate[i]] * 3; l <= (1200 - 3 * k) / 2; l++){
                if(dp[i - 1][9][k][l] < 0) continue;
                dp[i][9][k - wa[daystate[i]] * 3][l - fo[daystate[i]] * 3] = 
                        max(dp[i - 1][9][k][l] + 200, dp[i][9][k - wa[daystate[i]] * 3][l - fo[daystate[i]] * 3]);
                        //挖矿
            }

        // for(int k = 0; k <= 400; k++)   
        //     for(int l = 0; l <= (1200 - 3 * k) / 2; l++){
        //         if(k > 0)
        //             dp[i][39][k][l] = max(dp[i][39][k][l], dp[i][39][k - 1][l] - 10); //价格双倍
        //         if(l > 0)
        //             dp[i][39][k][l] = max(dp[i][39][k][l], dp[i][39][k][l - 1] - 20);
        //     }//买东西
    }

        double res = 0;
        int t = 0;
        int k1 = 0, l1 = 0;
        for(int i = 1; i <= 10; i++)
            for(int k = 0; k <= 400; k++)
                for(int l = 0; l <= 600; l++)
                    if(res < dp[i][13][k][l]){
                        res = dp[i][13][k][l];
                        k1 = k;
                        l1 = l;
                        t = i;
                    }
        dfs(t, 13, k1, l1, day);
        if(t == day) break;
    }
    
    for(int i = 0; i < cnt; i++){
        if(i == cnt - 1){
            st[i].re += st[i].wat * 2.5 + st[i].foo * 5.0;
            st[i].wat = st[i].foo = 0;
        }

        cout << st[i].t << " " << st[i].day << " " << st[i].wat << " " << st[i].foo << " " << st[i].re << endl;
    }
    for(int i = 0; i <= 10; i++)
        cout << daystate[i] << " ";
    cout << endl;
    auto end = std::chrono::high_resolution_clock::now();

    std::chrono::duration<double> duration = end - start;

    std::cout << "代码段运行时间: " << duration.count() << " 秒" << std::endl;
}
\end{minted}

\textbf{第四关代码}
\begin{minted}{c++}
#pragma GCC optimize("O2")
#pragma GCC optimize("Ofast")
#pragma GCC optimize("inline")
#pragma GCC optimize("omit-frame-pointer")
#pragma GCC optimize("unroll-loops")
#include<bits/stdc++.h>
using namespace std;
int dp[32][30][403][603]; // i, j, k, l, 时间，地点，水量，食物。
int daystate[35] = {0, 2, 2, 2, 2, 2, 2, 2, 2, 2, 2,
                        2, 2, 2, 2, 2, 2, 2, 2, 2, 2,
                        2, 2, 2, 2, 2, 2, 2, 2, 2, 2}; // 1, 2, 3晴朗，高温，沙暴；
int wa[5] = {0, 3, 9, 10};
int fo[5] = {0, 4, 9, 10};
vector<int> v[70];
int comes = 1, wt, fd;
void sta(){

    //第四关数据
    for(int i = 1; i <= 4; i++)  
        for(int j = 0; j < 4; j++){
            v[j * 5 + i].push_back(j * 5 + i + 1);
            v[j * 5 + i + 1].push_back(j * 5 + i);
            v[j * 5 + i].push_back(j * 5 + i + 5);
            v[j * 5 + i + 5].push_back(j * 5 + i);
        }
    for(int i = 1; i <= 4; i++){
        v[i * 5].push_back(i * 5 + 5);
        v[i * 5 + 5].push_back(i * 5);
    }
    for(int i = 1; i <= 4; i++){
        v[20 + i].push_back(21 + i);
        v[21 + i].push_back(20 + i);
    }
}

struct S{
    int t, day, wat, foo, re;
}st[40];
int cnt = 0;

void dfs(int t, int cur, int w, int f, int now){
    if(now == t){
        wt = w;
        fd = f;
        comes = cur;
        st[cnt++] = {t, cur, w, f, dp[t][cur][w][f]};
        return;
    }


    if(t == 0){
        st[cnt++] = {t, cur, w, f, dp[t][cur][w][f]};
        // printf("%d %d %d %d %d\n", t, cur, w, f, dp[t][cur][w][f]);
        return;
    }
    if(daystate[t] != 3){
        if(cur != 14){ //移动
            for(auto p : v[cur])
                if(dp[t - 1][p][w + 2 * wa[ daystate[t] ]][f + 2 * fo[ daystate[t] ]] == dp[t][cur][w][f]){
                    dfs(t - 1, p, w + 2 * wa[ daystate[t]], f + 2 * fo[ daystate[t]], now);
                    // printf("%d %d %d %d %d\n", t, cur, w, f, dp[t][cur][w][f]);
                    return;
                }
        }else{
            for(auto p : v[cur])
                for(int k = 2 * wa[daystate[t]]; k <= w + 2 * wa[daystate[t]]; k++) // 要走一步
                    for(int l = 2 * fo[daystate[t]]; l <= f + 2 * fo[daystate[t]]; l++)
                        if(dp[t - 1][p][k][l] == dp[t][cur][w][f] + 
                            10 * (w + 2 * wa[daystate[t]] - k) + 20 * (f + 2 * fo[daystate[t]] - l)){
                                dfs(t - 1, p, k, l, now);
                                // printf("%d %d %d %d %d\n", t, cur, w, f, dp[t][cur][w][f]);
                                return;
                            }
                        }
    }
    if(cur == 18)
        if(dp[t - 1][cur][w + 3 * wa[daystate[t]]][f + 3 * fo[daystate[t]]] == dp[t][cur][w][f] - 1000){
            dfs(t - 1, cur, w + 3 * wa[daystate[t]] , f + 3 * fo[daystate[t]], now);
            // printf("%d %d %d %d %d\n", t, cur, w, f, dp[t][cur][w][f]);
            return;
        }

    if(dp[t - 1][cur][w + 1 * wa[daystate[t]]][f + 1 * fo[daystate[t]]] == dp[t][cur][w][f]){
        dfs(t - 1, cur, w + 1 * wa[daystate[t]] , f + 1 * fo[daystate[t]], now);
        // printf("%d %d %d %d %d\n", t, cur, w, f, dp[t][cur][w][f]);
        return;
    }
    
}

int main(){
    ios::sync_with_stdio(false);
    cin.tie(0);
    auto start = std::chrono::high_resolution_clock::now();
    sta();
    

    double trans_prob[3][3];
    trans_prob[1][1] = 0.5;   // 晴→晴: 2次
    trans_prob[1][2] = 0.4;   // 晴→高: 6次
    trans_prob[1][3] = 0.1;        // 晴→沙: 0
        
        // 高温(2)的转移（只考虑晴朗和高温）
    trans_prob[2][1] = 0.3;  // 高→晴: 5次
    trans_prob[2][2] = 0.6;  // 高→高: 10次
    trans_prob[2][3] = 0.1;        // 高→沙: 0
        
        // 沙暴(3)的转移（实际不会用到）
    trans_prob[3][1] = 0.6;
    trans_prob[3][2] = 0.3;
    trans_prob[3][3] = 0.1;
    srand(time(0));
    
    for(int day = 1; day <= 30; day++){

    double ran = rand() % 50000 / (double)50000;
    if(daystate[day - 1] == 0){
        if(trans_prob[1][1] >= ran) {
            
            daystate[day] = 1; // 晴天
        }else if(trans_prob[1][1] + trans_prob[1][2] >= ran)
            daystate[day] = 2;// 高温
        else 
            daystate[day] = 3; //沙暴
    }else {
        if(trans_prob[daystate[day - 1]][1] >= ran) // 根据上一天的天气来看
            daystate[day] = 1; // 晴天
        else if(trans_prob[daystate[day - 1]][1] + trans_prob[daystate[day - 1]][2] >= ran)
            daystate[day] = 2;// 高温
        else 
            daystate[day] = 3; //沙暴
    }
    // daystate[day] = 2;
    for(int i = 0; i <= 30; i++)
        for(int j = 0; j <= 25; j++)
            for(int k = 0; k <= 400; k++)
                for(int l = 0; l <= 600; l++)
                    if(i != day - 1 || comes != j || day == 1 || wt != k || fd != l)
                        dp[i][j][k][l] = -5e3;
    if(day == 1){
        for(int k = 0; k <= 400; k++){
            for(int l = 0; l <= (1200 - 3 * k) / 2; l++){
                if(10000 - 5 * k - 10 * l >= 0)
                    dp[0][1][k][l] = 10000 - 5 * k - 10 * l;
                
            }
        }
    }

    for(int i = 1; i <= 30; i++){
        if(daystate[i] != 3) //非沙暴天
            for(int j = 1; j <= 25; j++)
                for(int k = wa[daystate[i]] * 2; k <= 400; k++)   
                    for(int l = fo[daystate[i]] * 2; l <= (1200 - 3 * k) / 2; l++){
                        if(dp[i - 1][j][k][l] < 0) continue;
                        for(auto t : v[j])
                            dp[i][t][k - wa[daystate[i]] * 2][l - fo[daystate[i]] * 2] = 
                            max(dp[i - 1][j][k][l], dp[i][t][k - wa[daystate[i]] * 2][l - fo[daystate[i]] * 2]);
                            //移动
                    }
        
        for(int j = 1; j <= 25; j++)
            for(int k = wa[daystate[i]]; k <= 400; k++)   
                for(int l = fo[daystate[i]]; l <= (1200 - 3 * k) / 2; l++){
                    if(dp[i - 1][j][k][l] < 0) continue;
                    dp[i][j][k - wa[daystate[i]]][l - fo[daystate[i]]] = 
                    max(dp[i - 1][j][k][l], dp[i][j][k - wa[daystate[i]]][l - fo[daystate[i]]]);
                            //停留
                }

        
        for(int k = wa[daystate[i]] * 3; k <= 400; k++)   
            for(int l = fo[daystate[i]] * 3; l <= (1200 - 3 * k) / 2; l++){
                if(dp[i - 1][18][k][l] < 0) continue;
                dp[i][18][k - wa[daystate[i]] * 3][l - fo[daystate[i]] * 3] = 
                        max(dp[i - 1][18][k][l] + 1000, dp[i][18][k - wa[daystate[i]] * 3][l - fo[daystate[i]] * 3]);
                        //挖矿
            }

        for(int k = 0; k <= 400; k++)   
            for(int l = 0; l <= (1200 - 3 * k) / 2; l++){
                if(k > 0)
                    dp[i][14][k][l] = max(dp[i][14][k][l], dp[i][14][k - 1][l] - 10); //价格双倍
                if(l > 0)
                    dp[i][14][k][l] = max(dp[i][14][k][l], dp[i][14][k][l - 1] - 20);
            }//买东西
    }

        double res = 0;
        int t = 0;
        int k1 = 0, l1 = 0;
        for(int i = 1; i <= 25; i++)
            for(int k = 0; k <= 400; k++)
                for(int l = 0; l <= 600; l++)
                    if(res < dp[i][25][k][l]){
                        res = dp[i][25][k][l];
                        k1 = k;
                        l1 = l;
                        t = i;
                    }
        
        dfs(t, 25, k1, l1, day);
        if(t == day) break;
    }
    
    for(int i = 0; i < cnt; i++){
        if(i == cnt){
            st[i].re += st[i].wat * 2.5 + st[i].foo * 5.0;
            st[i].wat = st[i].foo = 0;
        }

        cout << st[i].t << " " << st[i].day << " " << st[i].wat << " " << st[i].foo << " " << st[i].re << endl;
    }
    for(int i = 0; i <= 30; i++)
        cout << daystate[i] << " ";
    
    auto end = std::chrono::high_resolution_clock::now();

    std::chrono::duration<double> duration = end - start;

    std::cout << "代码段运行时间: " << duration.count() << " 秒" << std::endl;
}
\end{minted}

\textbf{第六关代码}
\begin{minted}{matlab}
function result = desert_survival_no_supply()
clc; clearvars -except result;

%% =========================
%  参数设置（可按需要修改）
%  =========================
P.deadline   = 30;      % 截止日期
P.moveDays   = 8;       % 起点到终点所需"可移动日"数
P.maxLoad    = 1200;    % 负重上限 kg
P.initMoney  = 10000;   % 初始资金 元

% 资源参数
P.waterMass  = 3;       % 每箱水质量 kg
P.foodMass   = 2;       % 每箱食物质量 kg
P.waterPrice = 5;       % 起点基准价格 元/箱
P.foodPrice  = 10;      % 起点基准价格 元/箱

%% ===== 挖矿问题新增参数 =====
P.villageWaterPrice = 2 * P.waterPrice;   % 村庄水价
P.villageFoodPrice  = 2 * P.foodPrice;    % 村庄食物价

P.distStartMine   = 5;   % 起点->矿山
P.distMineVillage = 2;   % 矿山->村庄（单程）
P.distMineEnd     = 3;   % 矿山->终点
P.distVillageEnd  = 3;   % 村庄->终点

P.mineIncome = 1000;     % 挖矿1天收益

% 补给次数序列：0 表示不去村庄；0.5表示最后只去一次村庄后直接去终点；
% 1 表示矿山->村庄->矿山 完整一来回；
% 1.5 表示一来回后，最后再去一次村庄并从村庄去终点
P.supplyCountList = 0:0.5:5;    % 这个范围你可以自己改

% 基础消耗（箱）
% [水, 食物]
P.baseClear  = [3, 4];
P.baseHot    = [9, 9];
P.baseSand   = [10, 10];

% 行走倍数、停留倍数（题目规则）
P.walkMult   = 2;       % 行走日 = 2 * 基础消耗
P.stayMult   = 1;       % 停留日 = 1 * 基础消耗

% 天气生成概率（这是"未知天气"下必须额外假设的部分，可自行改）
% 概率之和必须为1

P.pSand      = 0.15;
P.pClear     = 0.40;
P.pHot       = 1-P.pClear-P.pSand;

% 额外限制（按你截图中的思路）
P.maxSandDays = 10;     % c <= 10
P.maxTry      = 2e5;    % 为了采样足够有效路径，最多尝试次数
P.Nsim        = 5000;  % 蒙特卡洛样本数

%% =========================
%  计算"补给次数-最大盈利"关系
%  =========================
fprintf('正在计算"补给次数-最大盈利"关系...\n');

profitList = nan(size(P.supplyCountList));
bestMineDaysList = nan(size(P.supplyCountList));
bestInitWaterList = nan(size(P.supplyCountList));
bestInitFoodList  = nan(size(P.supplyCountList));

for k = 1:length(P.supplyCountList)
    supplyCount = P.supplyCountList(k);
    fprintf('正在计算补给次数 = %.1f ...\n', supplyCount);

    [bestProfit, bestMineDays, bestW0, bestF0] = ...
        optimizeMiningPlanForSupplyCount(P, supplyCount);

    profitList(k) = bestProfit;
    bestMineDaysList(k) = bestMineDays;
    bestInitWaterList(k) = bestW0;
    bestInitFoodList(k)  = bestF0;
end

%% 绘图
figure;
plot(P.supplyCountList, profitList, '-o', 'LineWidth', 1.5, 'MarkerSize', 6);
xlabel('去村庄补给次数');
ylabel('最大盈利');
title('去村庄补给次数与最大盈利的关系');
grid on;

%% 如需同时看看最优挖矿天数，可加第二张图
figure;
plot(P.supplyCountList, bestMineDaysList, '-s', 'LineWidth', 1.5, 'MarkerSize', 6);
xlabel('去村庄补给次数');
ylabel('最优挖矿天数');
title('去村庄补给次数与最优挖矿天数的关系');
grid on;

%% 输出结果
result.supplyCountList = P.supplyCountList;
result.profitList = profitList;
result.bestMineDaysList = bestMineDaysList;
result.bestInitWaterList = bestInitWaterList;
result.bestInitFoodList = bestInitFoodList;

%% 输出结果
result.sandProbList = sandProbList;
result.survivalRateList = survivalRateList;
result.bestWaterList = bestWaterList;
result.bestFoodList = bestFoodList;
end


%% =========================================================
%  蒙特卡洛生成样本：每条样本表示一条"完整到达终点"的天气路径
%  输出：
%    needW(k), needF(k): 第k条样本到达终点所需的水/食物箱数
% =========================================================
function [needW, needF, tripDays] = simulateDemands(P)

needW    = zeros(P.Nsim,1);
needF    = zeros(P.Nsim,1);
tripDays = zeros(P.Nsim,1);

cnt = 0;      % 已接受样本数
tries = 0;    % 总尝试次数

while cnt < P.Nsim && tries < P.maxTry
    tries = tries + 1;

    moved = 0;      % 已完成的可移动天数
    day   = 0;      % 已经历总天数
    sandDays = 0;

    seq = [];       % 1=晴朗, 2=高温, 3=沙暴
    w = 0; f = 0;

    valid = true;

    while day < P.deadline && moved < P.moveDays
        day = day + 1;
        wt = drawWeather(P);

        if wt == 3
            % 沙暴：必须停留
            sandDays = sandDays + 1;
            w = w + P.stayMult * P.baseSand(1);
            f = f + P.stayMult * P.baseSand(2);
        elseif wt == 1
            % 晴朗：前进
            moved = moved + 1;
            w = w + P.walkMult * P.baseClear(1);
            f = f + P.walkMult * P.baseClear(2);
        else
            % 高温：前进
            moved = moved + 1;
            w = w + P.walkMult * P.baseHot(1);
            f = f + P.walkMult * P.baseHot(2);
        end

        seq(end+1) = wt; %#ok<AGROW>

        % 约束1：沙暴总数上限
        if sandDays > P.maxSandDays
            valid = false;
            break;
        end

        % 约束2：到当前为止沙暴比例 < 1/3（你图里的"较少出现沙暴"）
        if sandDays / day >= 1/3
            valid = false;
            break;
        end

        % 约束3：任意连续 j>=2 天，沙暴比例 <= 1/2
        if ~checkWindowConstraint(seq)
            valid = false;
            break;
        end
    end

    % 必须在截止日前完成 moveDays 次有效移动
    if valid && moved >= P.moveDays && day <= P.deadline
        cnt = cnt + 1;
        needW(cnt) = w;
        needF(cnt) = f;
        tripDays(cnt) = day;
    end
end

% 截断未使用部分
needW    = needW(1:cnt);
needF    = needF(1:cnt);
tripDays = tripDays(1:cnt);

end


%% =========================================================
%  搜索最优初始购买方案
%  生存条件：初始水 >= 该样本所需水，且初始食物 >= 该样本所需食物
% =========================================================
function [bestW0, bestF0, bestRate] = searchBestPlan(needW, needF, P)

maxWaterByLoad  = floor(P.maxLoad / P.waterMass);
maxFoodByLoad   = floor(P.maxLoad / P.foodMass);
maxWaterByMoney = floor(P.initMoney / P.waterPrice);
maxFoodByMoney  = floor(P.initMoney / P.foodPrice);

maxW = min(maxWaterByLoad,  maxWaterByMoney);
maxF = min(maxFoodByLoad,   maxFoodByMoney);

bestRate = -1;
bestW0 = 0;
bestF0 = 0;

for W0 = 0:maxW
    % 由资金、负重双重限制推 F0 上界
    maxF_load  = floor((P.maxLoad  - P.waterMass  * W0) / P.foodMass);
    maxF_money = floor((P.initMoney - P.waterPrice * W0) / P.foodPrice);
    upperF = min([maxF, maxF_load, maxF_money]);

    if upperF < 0
        continue;
    end

    for F0 = 0:upperF
        alive = (W0 >= needW) & (F0 >= needF);
        rate = mean(alive);

        if rate > bestRate
            bestRate = rate;
            bestW0 = W0;
            bestF0 = F0;
        end
    end
end

end


%% =========================================================
%  随机生成一天的天气
%  1=晴朗, 2=高温, 3=沙暴
% =========================================================
function wt = drawWeather(P)
r = rand;
if r < P.pClear
    wt = 1;
elseif r < P.pClear + P.pHot
    wt = 2;
else
    wt = 3;
end
end


%% =========================================================
%  检查：任意连续 j>=2 天，沙暴比例 <= 1/2
%  这里采用严格判定
% =========================================================
function ok = checkWindowConstraint(seq)
n = length(seq);
ok = true;

for len = 2:n
    for s = 1:(n-len+1)
        window = seq(s:s+len-1);
        sandRatio = sum(window == 3) / len;
        if sandRatio > 1/2
            ok = false;
            return;
        end
    end
end
end

function [bestProfit, bestMineDays, bestW0, bestF0] = optimizeMiningPlanForSupplyCount(P, supplyCount)

bestProfit = -inf;
bestMineDays = NaN;
bestW0 = NaN;
bestF0 = NaN;

% 根据补给次数，先确定固定路线所需移动天数
[travelDays, needFinalVillageToEnd] = getTravelDaysFromSupplyCount(P, supplyCount);

% 超过截止日期则直接不可行
if travelDays > P.deadline
    return;
end

% 能用于挖矿的最大天数
maxMineDays = P.deadline - travelDays;

% 初始可购买上限（受资金和负重限制）
maxWaterByLoad  = floor(P.maxLoad / P.waterMass);
maxFoodByLoad   = floor(P.maxLoad / P.foodMass);
maxWaterByMoney = floor(P.initMoney / P.waterPrice);
maxFoodByMoney  = floor(P.initMoney / P.foodPrice);

maxW = min(maxWaterByLoad, maxWaterByMoney);
maxF = min(maxFoodByLoad, maxFoodByMoney);

% 枚举：初始水、初始食物、挖矿天数
for W0 = 0:maxW
    maxF_load  = floor((P.maxLoad - P.waterMass * W0) / P.foodMass);
    maxF_money = floor((P.initMoney - P.waterPrice * W0) / P.foodPrice);
    upperF = min(maxF, min(maxF_load, maxF_money));

    if upperF < 0
        continue;
    end

    for F0 = 0:upperF
        initCost = W0 * P.waterPrice + F0 * P.foodPrice;

        % 初始购买就已经超预算
        if initCost > P.initMoney
            continue;
        end

        for mineDays = 0:maxMineDays
            ok = checkMiningPlanFeasible(W0, F0, mineDays, supplyCount, P);

            if ok
                profit = mineDays * P.mineIncome - initCost;

                if profit > bestProfit
                    bestProfit = profit;
                    bestMineDays = mineDays;
                    bestW0 = W0;
                    bestF0 = F0;
                end
            end
        end
    end
end

if isinf(bestProfit)
    bestProfit = NaN;
end

end

function [travelDays, needFinalVillageToEnd] = getTravelDaysFromSupplyCount(P, supplyCount)

% supplyCount:
% 0   = 起点->矿山->终点
% 0.5 = 起点->矿山->村庄->终点
% 1   = 起点->矿山->村庄->矿山->终点
% 1.5 = 起点->矿山->村庄->矿山->村庄->终点
% 2   = 起点->矿山->村庄->矿山->村庄->矿山->终点
% ...

nRound = floor(supplyCount);              % 完整往返次数（矿山->村庄->矿山）
hasHalf = abs(supplyCount - nRound - 0.5) < 1e-9;

travelDays = P.distStartMine ...
           + nRound * (2 * P.distMineVillage);

if hasHalf
    % 最后一次只到村庄，再从村庄到终点
    travelDays = travelDays + P.distMineVillage + P.distVillageEnd;
    needFinalVillageToEnd = true;
else
    % 直接从矿山到终点
    travelDays = travelDays + P.distMineEnd;
    needFinalVillageToEnd = false;
end

end

function ok = checkMiningPlanFeasible(W0, F0, mineDays, supplyCount, P)

ok = false;

% 用高温作为保守消耗标准
walkNeed = P.walkMult * P.baseHot;   % [水, 食物]
mineNeed = 3 * P.baseHot;            % 挖矿1天 [水, 食物]

nRound = floor(supplyCount);
hasHalf = abs(supplyCount - nRound - 0.5) < 1e-9;

% -------- 第1段：起点 -> 矿山 --------
need1 = P.distStartMine * walkNeed;
W = W0 - need1(1);
F = F0 - need1(2);

if W < 0 || F < 0
    return;
end

% -------- 中间若干轮：矿山挖矿 + 去村庄补给 + 返回矿山 --------
% 设总挖矿天数尽量平均分布到各次在矿山阶段
numMineStages = nRound + 1;   % 完整往返n次时，共有 n+1 个"在矿山可挖矿阶段"
if hasHalf
    numMineStages = nRound + 1; % 最后半次前也仍有 n+1 个矿山阶段
end

mineDaysAlloc = floor(mineDays / numMineStages) * ones(1, numMineStages);
r = mod(mineDays, numMineStages);
mineDaysAlloc(1:r) = mineDaysAlloc(1:r) + 1;

for i = 1:nRound
    % 先在矿山挖矿
    W = W - mineDaysAlloc(i) * mineNeed(1);
    F = F - mineDaysAlloc(i) * mineNeed(2);
    if W < 0 || F < 0
        return;
    end

    % 从矿山走到村庄
    goNeed = P.distMineVillage * walkNeed;
    W = W - goNeed(1);
    F = F - goNeed(2);
    if W < 0 || F < 0
        return;
    end

    % 在村庄补给：为了完成"回矿山 + 后续所有行程"，一次性补到刚好够且不超负重
    [futureW, futureF] = futureNeedAfterVillage(i, nRound, hasHalf, mineDaysAlloc, P, walkNeed, mineNeed);

    buyW = max(0, futureW - W);
    buyF = max(0, futureF - F);

    if buyW * P.waterMass + buyF * P.foodMass > P.maxLoad
        return;
    end

    W = W + buyW;
    F = F + buyF;

    % 回矿山
    backNeed = P.distMineVillage * walkNeed;
    W = W - backNeed(1);
    F = F - backNeed(2);
    if W < 0 || F < 0
        return;
    end
end

% -------- 最后一个矿山阶段 --------
W = W - mineDaysAlloc(end) * mineNeed(1);
F = F - mineDaysAlloc(end) * mineNeed(2);
if W < 0 || F < 0
    return;
end

% -------- 结束段：去终点 --------
if hasHalf
    % 矿山->村庄
    goNeed = P.distMineVillage * walkNeed;
    W = W - goNeed(1);
    F = F - goNeed(2);
    if W < 0 || F < 0
        return;
    end

    % 最后一次只到村庄，补给后直接去终点
    endNeed = P.distVillageEnd * walkNeed;
    buyW = max(0, endNeed(1) - W);
    buyF = max(0, endNeed(2) - F);

    if buyW * P.waterMass + buyF * P.foodMass > P.maxLoad
        return;
    end

    W = W + buyW;
    F = F + buyF;

    W = W - endNeed(1);
    F = F - endNeed(2);
    if W < 0 || F < 0
        return;
    end
else
    % 直接矿山->终点
    endNeed = P.distMineEnd * walkNeed;
    W = W - endNeed(1);
    F = F - endNeed(2);
    if W < 0 || F < 0
        return;
    end
end

ok = true;

end

function [futureW, futureF] = futureNeedAfterVillage(iRound, nRound, hasHalf, mineDaysAlloc, P, walkNeed, mineNeed)

% 站在"第 iRound 次到达村庄补给点"这个时刻，
% 计算从此刻开始到终点还至少需要多少水/食物

futureW = 0;
futureF = 0;

% 先回矿山
futureW = futureW + P.distMineVillage * walkNeed(1);
futureF = futureF + P.distMineVillage * walkNeed(2);

% 后续矿山阶段：第 iRound+1, ..., 最后
for j = (iRound+1):length(mineDaysAlloc)
    futureW = futureW + mineDaysAlloc(j) * mineNeed(1);
    futureF = futureF + mineDaysAlloc(j) * mineNeed(2);

    if j < length(mineDaysAlloc)
        % 还要再去一次村庄并返回矿山
        futureW = futureW + 2 * P.distMineVillage * walkNeed(1);
        futureF = futureF + 2 * P.distMineVillage * walkNeed(2);
    end
end

% 最后一段去终点
if hasHalf
    futureW = futureW + P.distMineVillage * walkNeed(1) + P.distVillageEnd * walkNeed(1);
    futureF = futureF + P.distMineVillage * walkNeed(2) + P.distVillageEnd * walkNeed(2);
else
    futureW = futureW + P.distMineEnd * walkNeed(1);
    futureF = futureF + P.distMineEnd * walkNeed(2);
end

end
\end{minted}